\section{Introduction}

An idea that a ``type theory should eat itself'' is a
well-established concept in study of dependent type
theory. Under this guise, if a type theory is presented
as a ``full-scale system for constructive mathematics``
[CITE], then it should be capable of proving its own
metatheoretical properties. For example, in a paper
bearing the same name [CITE], the author presents a
normalizer for type theory written in Agda [CITE].

Of course, G\"odel's incompleteness theorems [CITE]
still apply in a constructive context, so na\"ively
proving self-consistency of a type system supporting
primitive recursion is out of question. What is usually
proven instead is that a particular implementation of a
type theory inside itself is sound assuming the
consistency of the theory, so no paradox arises [CITE]
[CITE]; this is still useful because it leads to the
phenomena of \emph{certified proof assistants}.

Another possible approach is to work in a two-level type
theory [CITE]. As the name implies, such a theory
consists of two levels: the lower \emph{object} level is
the main type theory, and a higher \emph{meta} level is
used to internalize statements about object level which
are external to it. A classical example given in the
cited paper is a construction of a \emph{semisimplicial
type of level $n$} for each external $n\in\mathbb{N}$;
object level here is homotopy type theory (HoTT) with
univalence axiom (UA) [CITE] and meta level is
Martin-L\"of type theory (MLTT) with uniqueness of
identity proofs (UIP) [CITE]; in usual type theory,
these two axioms are mutually exclusive, but the trick
is that an object-level equality becomes a different
type when viewed from meta level, so UA and UIP apply to
different notions of ``equality'', keeping the system
consistent (and sound with existing models of HoTT
[CITE]).

We aim to generalize this construction in the following
way. It can be said that an object level has a fixed
\emph{interpretation} in a meta-level; what if the
system allowed to vary the interpretation? In other
words, what if the system included a generic
interpretation operator denoted $m\llbracket t
\rrbracket$, where $t$ is a term which is being
interpreted, and $m$ is a model definable inside the
system? (We will call such systems \emph{reflective}.)

Turns out, we would get quite a few useful
metatheoretical results about object-level theory for
free; in particular, Wadler's parametricity theorem
[CITE] (for its formulation in dependently typed
context, see [CITE] and [CITE]) about object-level
theory would be provable on the meta-level by means of
specifying a relational model; a fraction of displayed
type theory [CITE] can be obtained via similar
translation of a type into a unary predicate. Another
fun consequence is that object-level theory would get
universe polymorphism via displacements as described by
Conor Titania McBride [CITE] and Favonia [CITE]; we
suspect that there might be other useful applications of
internal interpretation.

\subsection{Contributions}

So the goal is to construct type systems with
        such interpretation present as ; we will call type system with such
interpretation capabilities \emph{reflective}.
\begin{itemize}
    \item So our main goal is to construct an
        implementable reflective type system. The main
        difficulty in constructing such a system arises
        from a certain no-go result discussed in
        Section~\ref{no-go-section}, which forces us to
        limit the power of reinterpretation in order to
        keep type checking decidable.
    \item For ease of presentation, we first offer a
        reflective type system based on the
        Hindley-Milner type system, where object level
        is simply-typed lambda calculus with product
        types (STLC), called $\SSTLC$
        (Section~\ref{HM-R}). We prove that it is
        conservative with respect to STLC (hence
        strongly normalizing and sound), and that global
        type inference is decidable (section
        \ref{HM-R-properties}).
    \item Building upon the insights gathered from
        developing $\SSTLC$, we are finally equipped to
        dive into the reflective version of two-level
        type theory, where object level is vanilla MLTT,
        dubbed $\twoLTT$ (Section~\ref{MLTT-R}). We
        prove soundness and conservativity of it as well
        (Section~\ref{MLTT-R-properties}).
    \item As intended, by virtue of reinterpretation,
        $\twoLTT$ incorporates both a restricted form of
        universe polymorphism and external
        parametricity. And, of course, you can use it as
        a system for synthetic mathematics with internal
        reuse of synthetic results
        (Section~\ref{MLTT-R-applications}).
\end{itemize}
